\newpage
\section{Particle Physics}
\subsection{Definitions used in the field (of study):}
\newpage
\includepdf[height=0.8\pdfpageheight, pagecommand={\thispagestyle{fancy}}]{Physics/build/ParticleTimeline.pdf}
\newpage
\begin{itemize}
	\item dE/dx
	\item Transverse momentum \pT:\@\ momentum transverse to the beam axis.
	\item Transverse Energy \ET:\@\ Energy transverse to beam axis, deposited in calorimeter.
	\item Scalar Fields
	\item Classical Scattering
	\item Run 1: 2010--2012
	\item Run 2: 2015--2018
	\item Run 3: 2022--2026
	\item Cutflow: Volume of phase space (number of events) after each cut is applied. Used to visualize the cut restrictions as they are applied. 
	\item Validation Region / Control Region: A selection of data space that very close to signal selection region. Used to `validate' the background region. Usually just has same cuts as signal region with exception of one cut inverted.
	\item Quantization
	\item On-shell: 
	\item Branching Ratio: The fraction of particles that decay via a particular decay mode.
	\item Cross Section: A measure of the probability that a specific process will occur when particles
	      collide. It is typically measured in units of area, such as barns (1 barn = $\rm 10^{-28} m^2$ = $\rm 10^{-24} cm^2$).
	\item Luminosity: A measure of the number of particles per unit area per unit time in a collider. It is used to calculate the event rate for a given process.
	\item Integrated Luminosity: The total luminosity accumulated over a specific period of time. It is often expressed in inverse femtobarns (\(\rm fb^{-1}\)).
	\item Event Topologies:
\end{itemize}
\subsection{Width}
When particles decay immeasurably fast (\W, \Z, \h, t), width $\Gamma$ is used instead of mean lifetime $\tau$. 
\subsection{Higgs}
The Higgs boson \h\ explains how the vector bosons get their mass (why it's \W, \Z and $\gamma$) instead of all or none having mass. 
The \h\ also explains how fermions get their mass, and there is a correspoondance between Newtons first law. Interacting with the higgs field makes it difficult for a body to accelerate, which is analogous to inertia.
\subsection{Gamma Matrix Identities}
	\begin{minipage}{0.45\textwidth}
		\begin{align*}
			\gamma^{0\dagger}&=\gamma^0
			\\
			\gamma^{j\dagger}&=-\gamma^j\quad\qty(j=1,2,3)
			\\
			\gamma^0\gamma^{\mu\dagger}\gamma^0&=\gamma^\mu
			\\
			\gamma^{0}\gamma^{0}&=\mathbb{I}
			\\
			\gamma^\mu\gamma_\mu&=4
			\\
			\gamma^\mu\gamma_\nu\gamma_\mu&=-2\gamma_\nu
			\\
			\gamma^\mu a_\mu&=\gamma_\mu a^\mu=\slsh{a}
			\\
			\psi^\dagger\gamma^0&=\bar{\psi}
		\end{align*}
	\end{minipage}
	\begin{minipage}{0.48\textwidth}
		\begin{align*}
			\qty{\gamma_\mu,\gamma_\nu}&=\gamma_\mu\gamma_\nu+\gamma_\nu\gamma_\mu=2g_{\mu\nu}
			\\
			\gamma^\mu\gamma_\mu&=4
			\\
			\qty[\gamma^\mu,a_\mu]&=0
			\\
			\gamma^\mu a_\mu&=\gamma_\mu a^\mu=\slsh{a}
			\\
			\Tr(\mathbb{I})&=4
			\\
			\Tr(\gamma_\mu\gamma_\nu)&=\Tr(\gamma^\mu\gamma^\nu)=4
			\\
			\Tr(\gamma^\mu\gamma^\nu\gamma^\rho\gamma^\sigma)&=4(g^{\mu\nu}g^{\rho\sigma}-g^{\mu\rho}g^{\nu\sigma}+g^{\mu\sigma}g^{\nu\rho})
		\end{align*}
	\end{minipage}
	\\[0.25 cm]
	Gamma 5 matrix:
	\\
	\begin{minipage}{0.45\textwidth}
		\begin{align*}
			\gamma_5^\dagger&=\gamma_5
			\\
			\gamma_{5}\gamma_{5}&=\mathbb{I}
			\\
			\qty{\gamma_5,\gamma_\nu}&=0
			\\
			\Tr(\gamma_5)&=0
			\\
			\Tr(\gamma_\mu\gamma_5)&=0
			\\
			\Tr(\gamma_\mu\gamma_\nu\gamma_5)&=0
			\\
			\Tr(\gamma_\mu\gamma_\nu\gamma_\rho\gamma_5)&=0
			\\
			\Tr(\gamma_\mu\gamma_\nu\gamma_\rho\gamma_\sigma\gamma_5)&=4i\epsilon_{\mu\nu\rho\sigma}
		\end{align*}
	\end{minipage}
	\begin{minipage}{0.45\textwidth}
		\begin{align*}
			P_L&=\frac{1-\gamma_5}{2}
			\\
			P_R&=\frac{1+\gamma_5}{2}
			\\
			P_L\gamma_\mu&=\gamma_\mu P_R
			\\
			P_L(\gamma^\mu\gamma^\nu\dotsc\text{odd})P_L&=0
			\\
			P_L(\gamma^\mu\gamma^\nu\dotsc\text{even})P_L&=P_L\qty(\gamma^\mu\gamma^\nu\dotsc)
		\end{align*}
	\end{minipage}
	\\[0.25 cm]
	Quick Complex Conjugates:
	\\
	\begin{align*}
		\qty(\bar{u}\gamma^\mu v)^*&=\bar{v}\gamma^\mu u 
		\\
		\qty(\bar{u}\gamma^\mu P_Lv)^*&=\bar{v}\gamma^\mu P_L u
	\end{align*}
\newpage
\subsection{Mandelstam Variables}

\begin{tikzpicture}[baseline=(current bounding box.south)]% S CHANNEL Template diagram. 
    % Template diagrams are labelled from top to bottom, going left to right. 
    % p_{a,b,c,...} are incoming,
    % k_{1,2,3,...} are outgoing,
    % v_{1,2,3,...} are vertices. 
    \def\propAngle{0}% angle of virtual propagators (from horizontal)
    \def\pAngle{30}% angle of incoming particle (from horizontal)
    \def\kAngle{30}% angle of outgoing particle (from horizontal)
    \def\legLength{1.5}% length of external leg
    \def\propLength{2}% length of propagator

    \begin{feynman}
        \vertex (v_1) at (0,0);
        \vertex (v_2) at ($(v_1)+(\propAngle:\propLength)$);

        \vertex (p_a) at ($(v_1)-(-\pAngle:\legLength)$);
        \vertex (p_b) at ($(v_1)-( \pAngle:\legLength)$);
        \vertex (k_1) at ($(v_2)+( \kAngle:\legLength)$);
        \vertex (k_2) at ($(v_2)+(-\kAngle:\legLength)$);

        \draw[scalar, momentum=\((p_a+p_b)\)] (v_1) -- (v_2);
        
        \draw[scalar, momentum=\(p_a\)]       (p_a) -- (v_1);
        \draw[scalar, momentum'=\(p_b\)] (p_b) -- (v_1);
        \draw[scalar, momentum=\(k_1\)]       (v_2) -- (k_1);
        \draw[scalar, momentum'=\(k_2\)] (v_2) -- (k_2);
    \end{feynman}
	\node at ($(v_1)! 0.5!(v_2)-(0,\legLength+1)$) {\(s=(p_a+p_b)^2c^2=(k_1+k_2)^2c^2\)};% below, midway between v_1 and v_2
\end{tikzpicture}
~~
\begin{tikzpicture}[baseline=(current bounding box.south)]% T CHANNEL Template diagram. 
    % Template diagrams are labelled from top to bottom, going left to right. 
    % p_{a,b,c,...} are incoming,
    % k_{1,2,3,...} are outgoing,
    % v_{1,2,3,...} are vertices. 
    \def\propAngle{0}% angle of virtual propagators (from vertical)
    \def\pAngle{25}% angle of incoming particle (from horizontal)
    \def\kAngle{25}% angle of outgoing particle (from horizontal)
    \def\legLength{2}% length of external leg
    \def\propLength{1.5}% length of propagator
    
    \begin{feynman}
        \vertex (v_1) at (0,0);
        \vertex (v_2) at ($(v_1)+(\propAngle-90:\propLength)$);

        \vertex (p_a) at ($(v_1)-(-\pAngle:\legLength)$);
        \vertex (k_1) at ($(v_1)+( \kAngle:\legLength)$);
        \vertex (p_b) at ($(v_2)-( \pAngle:\legLength)$);
        \vertex (k_2) at ($(v_2)+(-\kAngle:\legLength)$);

        \draw[scalar, momentum=\(p_a\)] (p_a) -- (v_1);
        \draw[scalar, momentum=\(k_1\)] (v_1) -- (k_1);

        \draw[scalar, momentum=\((p_a-k_1)\)] (v_1) -- (v_2);
        
        \draw[scalar, momentum'=\(p_b\)] (p_b) -- (v_2);
        \draw[scalar, momentum'=\(k_2\)] (v_2) -- (k_2);
    \end{feynman}
	\node at ($(p_b)! 0.5!(k_2)-(0,1)$) {\(t=(p_a-k_1)^2c^2=(k_2-p_b)^2c^2\)};% below, midway between p_b and k_2
\end{tikzpicture}
~~
\begin{tikzpicture}[baseline=(current bounding box.south)]% U CHANNEL Template diagram. 
    % Template diagrams are labelled from top to bottom, going left to right.
    % p_{a,b,c,...} are incoming,
    % k_{1,2,3,...} are outgoing,
    % v_{1,2,3,...} are vertices. 
    \def\propAngle{0}% angle of virtual propagators (from vertical)
    \def\pAngle{25}% angle of incoming particle (from horizontal)
    \def\kAngle{10}% angle of outgoing particle (from horizontal)
    \def\legLength{2}% length of external leg
    \def\propLength{2}% length of propagator
    
    \begin{feynman}
        \vertex (v_1) at (0,0);
        \vertex (v_2) at ($(v_1)+(\propAngle-90:\propLength)$);

        \vertex (p_a) at ($(v_1)-(-\pAngle:\legLength)$);
        \vertex (k_1) at ($(v_1)+( \kAngle:\legLength+0.5)$);
        \vertex (p_b) at ($(v_2)-( \pAngle:\legLength)$);
        \vertex (k_2) at ($(v_2)+(-\kAngle:\legLength+0.5)$);

        \draw[scalar, momentum=\(p_a\)] (p_a) -- (v_1);
        \draw[scalar, momentum'=\(k_2\)] (v_1) -- (k_2);

        \draw[scalar, momentum'=\((p_a-k_2)\)] (v_1) -- (v_2);
        
        \draw[scalar, momentum'=\(p_b\)] (p_b) -- (v_2);
        \draw[scalar, momentum=\(k_1\)] (v_2) -- (k_1);
    \end{feynman}
	\node at ($(p_b)! 0.5!(k_2)-(0,1)$) {\(u=(p_a-k_2)^2c^2=(k_1-p_b)^2c^2\)};% below, midway between p_b and k_2
\end{tikzpicture}

\newpage
%\subsection{Feynman Rules}
%	\tikzfeynmanset{
%		every vertex={red},
%		every particle={blue},
%		every blob={draw=black, pattern color=black},
%		}
%	\begin{align*}
%		%Incoming anti/fermions
%		\feynmandiagram[small, baseline=(v1.base), horizontal=v1 to v2]
%		{
%			v1 [color=black, particle=$f$] -- [fermion] v2[blob]-- [draw=none] v3 [particle={\phantom{$f$}}, color=white]
%		};
%		&\quad\longleftrightarrow\quad
%		u(p,s)
%		\\
%		\feynmandiagram[small, baseline=(v1.base), horizontal=v1 to v2]
%		{
%			v1 [color=black, particle=$f$] -- [anti fermion] v2[blob]-- [draw=none] v3 [particle={\phantom{$f$}}, color=white]
%		};
%		&\quad\longleftrightarrow\quad
%		\overline{v}(p,s)
%		\\
%		% Outgoing anti/fermions
%		\feynmandiagram[small, baseline=(v1.base), horizontal=v1 to v2]
%		{
%			v3 [draw=none] -- [draw=none] v1 [blob] -- [fermion] v2 [color=black, particle=$f$]
%		};
%		&\quad\longleftrightarrow\quad
%		\overline{u}(p,s)
%		\\
%		\feynmandiagram[small, baseline=(v1.base), horizontal=v1 to v2]
%		{
%			v3 [draw=none] -- [draw=none] v1 [blob] -- [anti fermion] v2 [color=black, particle=$f$]
%		};
%		&\quad\longleftrightarrow\quad
%		v(p,s)
%		\\
%		% Photon fermion interaction
%		\feynmandiagram[small, baseline=(v1.base), horizontal=v1 to v2]{
%		v1 [particle=$\mu$, color=black] -- [photon] v2,
%		[particle=$f$] a -- [fermion] v2 -- [fermion] b [particle=$f$],
%		};
%		&\quad\longleftrightarrow\quad
%		-iQ_fe\gamma^{\mu}
%	\end{align*}%